\section{Hexagon bundles over graphs}\label{sec:graphs}

We associate a toric fibre bundle to a graph whose vertices carry positive
integers.  Each edge contributes a degree-six del Pezzo surface, and each
vertex contributes a projective-space factor to the base.  The graph records
which two base factors twist a given surface factor.

\begin{definition}\label{graph:construction}
Let $\Gamma$ be a finite nonempty graph without loops, with vertex set $C$
and edge set $I$.  Multiple edges are allowed.  Write $k=|C|$ and $j=|I|$,
choose an orientation $c_i^-\to c_i^+$ for every edge $i$, and assign an
integer $b_c\ge1$ to each vertex $c$.  In the lattice
\[
 N=\bigoplus_{i\in I}\ZZ^2_i\ \oplus\ \bigoplus_{c\in C}\ZZ^{b_c}
\]
write $r_i=(1,0)_i$, $q_i=(0,1)_i$, and let
$f^c_1,\ldots,f^c_{b_c}$ be the standard basis of the last summand.  Set
\begin{equation}\label{graph:base-rays}
 f^c_0=-\sum_{t=1}^{b_c}f^c_t+
       \sum_{c_i^+=c}r_i-\sum_{c_i^-=c}r_i.
\end{equation}
For each edge take the six rays generated by
$\pm r_i,\pm q_i,\pm(r_i+q_i)$.  Together with the rays generated by all
$f^c_t$, these are the rays of $\Sigma_\Gamma$.  Its maximal cones contain
two cyclically adjacent hexagon rays for each $i$, and all but one of
$f^c_0,\ldots,f^c_{b_c}$ for each $c$; all their faces are included.
We write $X(\Gamma,b)$, or $X(\Gamma)$ when the integers are understood,
for the associated toric variety.
\end{definition}

The construction indeed gives a smooth projective fan, as proved below.
The surface with the indicated hexagon fan is the blow-up of $\PP^2$ at
its three torus-fixed points; we denote it by $S_6$, with the subscript
indicating its anticanonical degree.  Projection onto the last summands
will identify $X(\Gamma)$ as a toric fibre bundle with fibre $S_6^j$ over
$\prod_c\PP^{b_c}$.  Reversing the orientation of an edge gives an
isomorphic fan: multiplication by $-1$ on its summand $\ZZ^2_i$, and the
identity on every other summand, gives the required lattice isomorphism.
Thus no choice of edge orientation affects the variety.

We first record the primitive collections needed to test ampleness.  A
\emph{primitive collection} of a simplicial fan is a set of rays which
does not span a cone, whereas every proper subset does.  For a complete
simplicial fan, its \emph{primitive relation} expresses the sum of these
primitive ray generators as a positive linear combination of the ray
generators of the unique cone whose relative interior contains that sum.
The zero sum has the empty cone on the right.

\begin{lemma}\label{graph:split-collections}
Suppose a complete simplicial fan $\Sigma$ is split by a base fan
$\Sigma_B$ and a fibre fan $\Sigma_F$: there is a lattice projection with
kernel $N_F$ and a subfan $\widehat\Sigma_B$ projecting cone by cone
isomorphically onto $\Sigma_B$, and the cones of $\Sigma$ are exactly
$\widehat\sigma+\tau$, with $\widehat\sigma\in\widehat\Sigma_B$ and
$\tau\in\Sigma_F$.  The primitive collections of $\Sigma$ are precisely
the fibre primitive collections and the lifts of the base primitive
collections.  Fibre primitive relations are unchanged.  In a lifted
base primitive relation, the coefficients of the lifted base rays on
the right are those of the base relation; the other terms are a
nonnegative combination of fibre rays.
\end{lemma}

\begin{proof}
The two summands of $\widehat\sigma+\tau$ have independent linear spans:
projection is injective on the first and vanishes on the second.
Consequently a set of rays spans a cone exactly when its fibre part
spans a fibre cone and its base part projects to a base cone.  The
simplicial complex of $\Sigma$ is therefore the join of the complexes of
$\Sigma_F$ and $\Sigma_B$.  A minimal nonface of a join is a minimal
nonface of one factor, proving the assertion about collections.

Let $x$ be the sum of the generators of a lifted base collection, and
write its minimal containing cone as $\widehat\sigma+\tau$.  Projection
maps the relative interior of this cone onto the relative interior of
$\sigma$.  Thus $\sigma$ is the minimal containing cone of the projected
sum.  Comparing coefficients in the linearly independent generators of
$\sigma$ identifies all base coefficients; the remaining coefficients
are the nonnegative fibre coefficients in $x$.  If instead $x\in N_{F,\RR}$,
its minimal containing cone has $\widehat\sigma=0$, so the fibre relation
is unchanged.
\end{proof}

\begin{proposition}\label{graph:geometry}
The variety $X(\Gamma)$ is smooth and projective, and is a Zariski locally
trivial toric fibre bundle
\[
 X(\Gamma)\longrightarrow\prod_{c\in C}\PP^{b_c}
 \qquad\text{with fibre }S_6^j.
\]
Its dimension and Picard number are
\begin{equation}\label{graph:dimension-rank}
 n=2j+\sum_c b_c,\qquad \rho=4j+k.
\end{equation}
Its ray matroid is connected if and only if $\Gamma$ is connected.
Furthermore,
\begin{equation}\label{graph:fano-condition}
 X(\Gamma)\text{ is Fano}\quad\Longleftrightarrow\quad
 \deg_\Gamma(c)\le b_c\quad\text{for every }c\in C,
\end{equation}
where the degree counts edges with multiplicity.
\end{proposition}

\begin{proof}
We verify completeness and projectivity directly.  Let
\begin{equation}\label{graph:dual-hexagon}
 Q=\{(a,z)\in\RR^2: |a|\le1,\ |z|\le1,\ |a+z|\le1\},
\end{equation}
the anticanonical polygon of $S_6$.  In coordinates
$(m_i,y^c_t)$ on $M_\RR$, with $m_i=(a_i,z_i)$, put
$s_c=\sum_{c_i^+=c}a_i-\sum_{c_i^-=c}a_i$.
Choose integers $A_c>\deg_\Gamma(c)$ and consider the polytope
\[
 m_i\in Q,\qquad y^c_t\ge-1,\qquad
 \sum_{t=1}^{b_c}(y^c_t+1)\le A_c+s_c.
\]
Every slice over $Q^j$ is a product of full-dimensional simplices,
since $A_c+s_c>0$.  At any point at most $b_c$ inequalities of the
$c$th simplex are active.  A vertex therefore requires $2j$ independent
active hexagon inequalities, so its projection is a vertex of $Q^j$.
Conversely, each such vertex and each choice of a vertex in every
simplex give a vertex of the full polytope.  Exactly $n$ inequalities
are active, and their inward normals are precisely the generators of a
maximal cone in Definition~\ref{graph:construction}.
The fan of inward normals is consequently $\Sigma_\Gamma$; in particular
it is complete and projective.  Each maximal cone is unimodular: its
hexagon generators give a basis of $\bigoplus_i\ZZ^2_i$, and modulo that
sublattice its remaining generators form a standard maximal cone of
$\prod_c\PP^{b_c}$.  This also shows that the polytope just constructed
has integral vertices.

Projection to the base splits the fan in the sense of
Lemma~\ref{graph:split-collections}, with fibre the product hexagon fan.
The toric local triviality theorem gives the asserted bundle
\cite[Theorem~3.3.19]{CoxLittleSchenck2011}.
There are $6j+\sum_c(b_c+1)$ rays.  Subtracting $n$ gives
\eqref{graph:dimension-rank}.

The fibre primitive collections are the nonadjacent pairs in a
hexagon.  Their relations are either $v+(-v)=0$, of degree $2$, or
$v+v''=v'$, of degree $1$.  Here the degree of a primitive relation
is the number of terms on the left minus the sum of coefficients on
the right.  The other primitive collections, by
Lemma~\ref{graph:split-collections}, are
$\{f^c_0,\ldots,f^c_{b_c}\}$, with relations
\begin{equation}\label{graph:primitive-relations}
 f^c_0+\cdots+f^c_{b_c}
 =\sum_{c_i^+=c}r_i+\sum_{c_i^-=c}(-r_i).
\end{equation}
The right-hand side consists of one ray from each incident hexagon,
so these rays span a cone and all coefficients are $1$.  Its degree
is $b_c+1-\deg_\Gamma(c)$.  Batyrev's primitive-collection criterion
for ampleness, in the form of
\cite[Theorem~6.4.9(b)]{CoxLittleSchenck2011}, says that $-K_X$ is ample
exactly when all these degrees are positive.  Since the degrees are
integers, this is \eqref{graph:fano-condition}.

For the matroid assertion, the six rays of each hexagon belong to one
matroid component: the antipodal pairs and the relation
$r_i+q_i-(r_i+q_i)=0$ connect all six.  For each $c$, the set
\[
 \{f^c_0,\ldots,f^c_{b_c}\}\ \cup\ \{r_i:i\text{ is incident to }c\}
\]
is a circuit.  Indeed, projection to $\ZZ^{b_c}$ forces the
coefficients of all the $f^c_t$ in any relation to agree; the independent
$r_i$ then determine all the other coefficients, and a nonzero relation
has no zero coefficient.  These circuits join exactly the vertex and
edge factors incident in $\Gamma$.  Thus a connected graph gives a
connected matroid.  Distinct connected components of $\Gamma$ use
disjoint lattice summands, giving a direct sum of nonempty ray
matroids and proving the converse.
\end{proof}

\begin{lemma}[Recovery of the graph from the fan]\label{graph:reconstruction}
Let $\Gamma$ be a connected loopless graph with at least one edge.
The fan $\Sigma_\Gamma$ of Definition~\ref{graph:construction}
determines $\Gamma$ and the integers $b_c$ up to isomorphism.
In particular, isomorphic toric fans give isomorphic graphs with the
same vertex dimensions.
\end{lemma}

\begin{proof}
We first recover the fibre rays without using the given presentation.
They are precisely the primitive ray generators whose negatives also
occur in the fan. Every fibre ray has this property. A base ray cannot
be opposite a fibre ray, since its projection to the base lattice is
nonzero. Rays belonging to different base factors cannot be opposite,
since their projections lie in independent summands. Within a base
factor of dimension at least two, the standard projective-space rays
contain no opposite pair. If $b_c=1$, its two lifted rays have sum
\[
 f^c_0+f^c_1=\sum_{c_i^+=c}r_i-\sum_{c_i^-=c}r_i.
\]
This sum is nonzero: connectedness and the existence of an edge exclude
an isolated vertex, and the vectors from different incident edges
belong to independent hexagon summands. The loopless hypothesis ensures
that no edge contributes both signs at $c$.

Restrict the ray matroid to the recovered fibre rays. Its connected
components are exactly the six-ray sets of the individual hexagons.
Indeed, the opposite pairs and the circuit on the three ray lines
connect all rays of one hexagon, whereas independent hexagon planes
admit no circuit meeting two factors. We therefore recover each plane
$W_i$ and their direct sum $W=\bigoplus_iW_i$.

The primitive collections consisting entirely of the remaining rays
are precisely the groups
$C_c=\{f^c_0,\ldots,f^c_{b_c}\}$, by
Lemma~\ref{graph:split-collections}. Thus these groups, and the integers
$b_c=|C_c|-1$, are determined by the fan. Their vector sums lie in $W$:
\[
 s_c:=\sum_{u\in C_c}u
     =\sum_{c_i^+=c}r_i-\sum_{c_i^-=c}r_i.
\]
The component of $s_c$ in the recovered direct summand $W_i$ is nonzero
exactly when $c$ is an endpoint of $i$. This recovers the incidence
graph, including parallel edges, and hence $\Gamma$. Primitive ray
generators, opposite pairs, matroid components, primitive collections,
and these vector sums are all preserved by a lattice isomorphism of
fans, proving the assertion.
\end{proof}

\begin{corollary}\label{graph:ceiling}
If $\Gamma$ is connected and $X(\Gamma)$ is Fano, then
\begin{equation}\label{graph:rank-bound}
 \rho\le5\left\lfloor\frac n4\right\rfloor+1,
 \qquad \rho-n\le\frac n4+1.
\end{equation}
Equality in the first bound holds exactly when $\Gamma$ is a tree
with $j=\lfloor n/4\rfloor$ edges.  In that case
\[
 \sum_c\bigl(b_c-\deg_\Gamma(c)\bigr)=n-4\lfloor n/4\rfloor.
\]
In particular, for $j\ge1$ and $n=4j$, equality is equivalent to a tree
with $b_c=\deg_\Gamma(c)$ at every vertex.
\end{corollary}

\begin{proof}
Connectedness gives $k\le j+1$, with equality exactly for a tree.
The Fano condition gives $\sum_c b_c\ge\sum_c\deg_\Gamma(c)=2j$.
Hence $n\ge4j$ and
\[
 \rho=4j+k\le5j+1\le5\lfloor n/4\rfloor+1.
\]
Equality in this chain holds precisely when $k=j+1$ and
$j=\lfloor n/4\rfloor$, giving the stated characterization.
The formula for the sum follows from
$n-4j=\sum_c(b_c-\deg_\Gamma(c))$.
\end{proof}

The corollary includes the graph with one vertex and no edges:
$X(\Gamma)=\PP^{b_c}$, $\rho=1$, and equality in its first bound holds
exactly for $1\le b_c\le3$.  The empty graph is excluded throughout.
If $\Gamma$ is disconnected, the lattice and fan decompose over its
connected components, so $X(\Gamma)$ is a product of positive-dimensional
varieties.  Its tangent bundle is not stable for any polarization:
it splits into nonzero direct summands, and at least one summand has
slope at least the slope of their direct sum.  The bound
\eqref{graph:rank-bound} concerns this graph construction; it makes no
claim for arbitrary smooth toric Fano varieties.

\begin{example}\label{graph:path}
For $j\ge1$, let $\Gamma_j$ be the path with vertices $0,\ldots,j$,
edges $i:i\to i+1$ for $0\le i<j$, and weights
$b_0=b_j=1$, $b_c=2$ for $0<c<j$.
We denote $X(\Gamma_j)$ by $X_j$.
It is Fano, has connected ray matroid, and satisfies
\[
 \dim X_j=4j,\qquad \rho(X_j)=5j+1.
\]
Thus these paths attain the bound in Corollary~\ref{graph:ceiling}.
Connectedness alone does not establish stability; the slope inequalities
for this family will require the degree formulas below.
\end{example}

\subsection{Anticanonical degrees}

Assume henceforth that $b_c\ge\deg_\Gamma(c)$.  For
$a=(a_i)\in[-1,1]^j$ define
\begin{equation}\label{graph:density}
 s_c(a)=\sum_{c_i^+=c}a_i-\sum_{c_i^-=c}a_i,
 \qquad L_c(a)=b_c+1+s_c(a),\qquad
 \Phi(a)=\prod_{c\in C}\frac{L_c(a)^{b_c}}{b_c!}.
\end{equation}
In particular $L_c\ge1$.  All lattice volume measures below assign
volume $1$ to a fundamental parallelepiped, without a factorial.
Put
\[
 d\nu(a)=\prod_{i\in I}(2-|a_i|)\,da_i.
\]
For a hexagon ray $v$ in the $i$th factor, let $d\nu^{(i,v)}$ be
obtained by replacing the $i$th factor of $d\nu$ with the measure
$dE_v$ on $[-1,1]$ specified by
\begin{equation}\label{graph:edge-measures}
 \begin{aligned}
 dE_{r_i}&=\delta_{-1},& dE_{-r_i}&=\delta_{1},\\
 dE_{r_i+q_i}=dE_{-q_i}&=\mathbf1_{[-1,0]}(a_i)\,da_i,
 &dE_{q_i}=dE_{-r_i-q_i}&=\mathbf1_{[0,1]}(a_i)\,da_i.
 \end{aligned}
\end{equation}
Here $\delta_{\pm1}$ denote unit point masses, not normalized divisor
degrees.  Each measure in \eqref{graph:edge-measures} has total mass $1$.

\begin{proposition}\label{graph:degree-formulas}
Write $D=(-K_X)^n$ and $d_v=D_v\cdot(-K_X)^{n-1}$ for $X=X(\Gamma)$.
Then
\begin{align}
 D&=n!\int_{[-1,1]^j}\Phi\,d\nu,\label{graph:volume-formula}\\
 d_{f^c_t}&=(n-1)!\int_{[-1,1]^j}\frac{b_c}{L_c}\Phi\,d\nu
       &&(0\le t\le b_c),\label{graph:base-degree}\\
 d_v&=(n-1)!\int_{[-1,1]^j}\Phi\,d\nu^{(i,v)}
       &&(v\text{ in hexagon }i).\label{graph:fibre-degree}
\end{align}
These formulas also hold for $j=0$, with the integral over a point.
\end{proposition}

\begin{proof}
The moment polytope of $-K_X$ is given by
\[
 m_i=(a_i,z_i)\in Q,\qquad y^c_t\ge-1,\qquad
 -\sum_{t=1}^{b_c}y^c_t+s_c(a)\ge-1.
\]
After translating $y^c_t$ by $1$, its slice over $(m_i)\in Q^j$ is
\[
 \prod_c\{(u^c_t):u^c_t\ge0,\ \sum_tu^c_t\le L_c(a)\},
\]
whose volume is $\Phi(a)$.  For fixed $a\in[-1,1]$, the allowable
$z$ in \eqref{graph:dual-hexagon} form an interval of length $2-|a|$.
Fubini's theorem and $D=n!\operatorname{vol}(P_{-K_X})$ therefore give
\eqref{graph:volume-formula}.

For $t\ge1$, the facet corresponding to $f^c_t$ sets $u^c_t=0$.
It replaces one simplex volume by
$L_c^{b_c-1}/(b_c-1)!$, multiplying the integrand by $b_c/L_c$.
The facet corresponding to $f^c_0$ sets
$\sum_tu^c_t=L_c(a)$; eliminating one coordinate gives the same
volume.  This elimination is an isomorphism of the affine facet
lattice with the coordinate lattice, since the coefficient of the
eliminated coordinate is $1$ and $s_c$ is integral affine.
Thus no additional factor occurs in its lattice measure.  The standard
facet-volume identity $d_v=(n-1)!\operatorname{vol}(F_v)$ proves
\eqref{graph:base-degree}, including $b_c=1$.

A facet for a hexagon ray lies over the corresponding edge of $Q$.
Each such edge has lattice length $1$.  The edges $a=-1$ and $a=1$
project to the unit point masses in \eqref{graph:edge-measures}.
On each other edge, projection to $a$ is a lattice isomorphism onto
an interval of length $1$.  Specifically, $a+z=-1$ and $z=1$ give
$[-1,0]$, whereas $z=-1$ and $a+z=1$ give $[0,1]$.  These are exactly
the four interval measures in \eqref{graph:edge-measures}.
Fubini on the facet proves \eqref{graph:fibre-degree}.
\end{proof}

\begin{remark}\label{graph:normalized-degrees}
If $Z=\int\Phi\,d\nu$, the normalized degrees
$\delta_v=nd_v/D$ are the respective integrals in
\eqref{graph:base-degree} and \eqref{graph:fibre-degree}, divided by $Z$.
Thus the factorials cancel before applying the slope criterion.
In particular the $b_c+1$ rays associated to a vertex have equal degrees.
\end{remark}
